% Part: first-order-logic
% Chapter: introduction
% Section: syntax

\documentclass[../../../include/open-logic-section]{subfiles}

\begin{document}

\olfileid[fr]{fol}{int}{syn}

\olsection{Syntaxe}

Nous devons d'abord préciser quelles chaînes de symboles sont des
!!{sentence}s de la logique du premier ordre. Nous le ferons plus loin ;
pour l'instant, procédons seulement par l'exemple. Les éléments de
base --- le vocabulaire --- de la logique du premier ordre se divisent
en deux parties. La première comprend les symboles qui servent à
désigner des choses particulières ou à en dire quelque chose. Nous
désignons des choses à l'aide de !!{constant}s, et nous parlons des
choses ainsi désignées à l'aide de !!{predicate}s. Nous pouvons par
exemple employer $\Obj a$ comme !!a{constant} pour désigner une chose
unique, puis en dire quelque chose au moyen de
l'!!{sentence}~$\Atom{\Obj P}{\Obj a}$. Si vous avez en tête une
signification pour « $\Obj a$ » et « $\Obj P$ », vous pouvez lire
$\Atom{\Obj P}{\Obj a}$ comme une phrase de la langue courante (et vous
l'avez probablement déjà fait lors de votre première initiation à la
logique formelle). À partir de tels !!{sentence}s simples de la logique
du premier ordre, on peut en construire de plus complexes à l'aide de
la seconde partie du vocabulaire : les symboles logiques, c'est-à-dire
les connecteurs et les quantificateurs. On peut ainsi former des
expressions comme $(\Atom{\Obj P}{\Obj a} \land
\Atom{\Obj Q}{\Obj b})$ ou~$\lexists[\Obj x][\Atom{\Obj P}{\Obj x}]$.

Pour donner les définitions précises de la sémantique et des règles de
nos systèmes de !!{derivation}, indispensables à une étude métalogique
rigoureuse, il faut d'abord définir exactement ce qui constitue
!!a{sentence} de la logique du premier ordre. L'idée générale est assez
simple. Certains !!{sentence}s élémentaires, comme
$\Atom{\Obj P}{\Obj a}$, ne font intervenir que des !!{predicate}s et
des !!{constant}s ; on en forme ensuite de plus complexes au moyen des
connecteurs et des quantificateurs. Mais quelles sont exactement les
règles qui autorisent la formation de ces !!{sentence}s plus
complexes ? Il faut les énoncer, faute de quoi l'expression
« !!{sentence} de la logique du premier ordre » ne serait pas définie
avec assez de précision. Plusieurs difficultés se présentent. Il faut
d'abord déterminer exactement quelles chaînes sont des !!{sentence}s. Il
faut ensuite le faire d'une manière qui permette de démontrer des
résultats mathématiques sur \emph{tous} les !!{sentence}s. Enfin, il
faudra définir précisément quelques opérations élémentaires sur les
!!{sentence}s, telles que « remplacer chaque $\Obj x$ de~$!A$
par~$\Obj b$ ». Les quantificateurs et les !!{variable}s de la logique
du premier ordre rendent la bonne définition de cette opération moins
évidente qu'il n'y paraît. Par exemple,
$\lexists[\Obj x][\Atom{\Obj P}{\Obj a}]$ doit-il être considéré comme
!!a{sentence} ? Qu'en est-il de
$\lexists[\Obj x][\lexists[\Obj x][\Atom{\Obj P}{\Obj x}]]$ ? Et quel
doit être le résultat de « remplacer $\Obj x$ par~$\Obj b$ dans
$(\Atom{\Obj P}{\Obj x} \land
\lexists[\Obj x][\Atom{\Obj P}{\Obj x}])$ » ?

\end{document}
